\section{Preliminaries}
\label{sec:prelim}

This section introduces the two mathematical tools the rest of the
report builds on: uniform B-splines as a continuous-time trajectory
representation, both in $\mathbb{R}^3$ and, in cumulative form, on
$\SO$; and the \ac{map} estimator that fits them to the measurements as
a weighted nonlinear least-squares problem. 

\subsection{Uniform B-splines in $\mathbb{R}^3$}
\label{sec:prelim_bspline}
An order-$N$ (degree $N{-}1$) uniform B-spline with control points
$\{\mathbf{P}_i\}\subset\mathbb{R}^3$ and constant knot spacing $\Delta t$
represents a curve as a weighted sum of basis functions,
\begin{equation}
  \vp(t) = \sum_i \mathbf{P}_i\, B_{i,N}(t),
  \label{eq:bspline_def}
\end{equation}
where the $B_{i,N}$ are the Cox--de Boor basis
functions~\cite{furgale2013unified}.  Two properties make this
representation attractive for asynchronous sensor fusion.

\emph{Local support.}  At any instant only $N$ consecutive control
points have nonzero basis, so on the segment $t\in[t_i,t_{i+1})$, with
normalized time $u=(t-t_i)/\Delta t\in[0,1)$, the curve is a fixed
blending of $N$ control points,
\begin{equation}
  \begin{gathered}
  \vp(t) = \sum_{j=0}^{N-1} B_j(u)\,\mathbf{P}_{i+j}, \\
  \bigl[B_0(u)\,\cdots\,B_{N-1}(u)\bigr]
     = \mathbf{u}^{\top}\mathbf{M}^{(N)},
  \end{gathered}
  \label{eq:bspline_blend}
\end{equation}
with $\mathbf{u}=[1,u,\dots,u^{N-1}]^{\top}$ and a constant $N\times N$
blending matrix $\mathbf{M}^{(N)}$~\cite{usenko2020basalt}.  Local
support limits each measurement's Jacobian to $N$ control points and
supports the marginalization described in \autoref{sec:sw}.

\emph{Analytic derivatives.}  Differentiating the basis in
\eqref{eq:bspline_blend} yields the velocity $\dot{\vp}(t)$ and
acceleration $\ddot{\vp}(t)$ in closed form at \emph{any} $t$, as
lower-order splines over the same control points, with no finite
differencing or inter-sample interpolation.  This is what lets the
accelerometer factor (which needs $\ddot{\vp}$) and the min-snap
regularizer (which needs $\vp^{(4)}$) be evaluated exactly at each
measurement's timestamp.  We use a quintic ($N{=}6$) position spline so
that snap is well defined and continuous.

\subsection{Cumulative $\SO$ B-splines}
\label{sec:prelim_so3}
Orientation lives on the manifold $\SO$, on which a weighted sum of
rotations $\sum_i \vR_i B_i$ is not itself a rotation, so
\eqref{eq:bspline_def} does not transfer directly.  The
\emph{cumulative} parameterization of Sommer et
al.~\cite{sommer2020efficient}, used in the basalt
\ac{vio}~\cite{usenko2020basalt}, instead writes the orientation as a
base rotation followed by a product of incremental rotations, each
mapped back onto the manifold by the exponential map,
\begin{equation}
  \begin{aligned}
  \vR(t) &= \vR_{i}
           \prod_{j=1}^{N-1}
           \Exp\!\bigl(\tilde{B}_j(u)\,\vOm_j\bigr), \\
  \vOm_j &= \Log\!\bigl(\vR_{i+j-1}^{-1}\vR_{i+j}\bigr)\in\so,
  \end{aligned}
  \label{eq:prelim_so3}
\end{equation}
where the increments $\vOm_j$ are the relative rotations between
consecutive \emph{absolute} orientation control points $\vR_i\in\SO$
(the optimization variables, basalt convention).  These orientation
control points are also called ``knots'' in the continuous-time
estimation literature~\cite{sommer2020efficient}; here, we use
``control points'' for both coefficient sets and reserve ``knots''
for the time values defining the spline intervals.  The cumulative
basis is the tail sum of the ordinary basis of \eqref{eq:bspline_blend},
\begin{equation}
  \tilde{B}_j(u) = \sum_{l=j}^{N-1} B_l(u).
  \label{eq:cumbasis}
\end{equation}
Because each $\Exp(\cdot)$ returns a rotation, $\vR(t)\in\SO$ holds by
construction for any orientation control points.  This form retains
absolute orientation control points, preserves local support over $N$
control points exactly as in the $\mathbb{R}^3$ case, and yields the body
angular velocity $\vom_b(t)$ in
closed form by differentiating the product, with no numerical
differentiation.  We use a cubic ($N{=}4$) orientation spline;
\eqref{eq:ori_spline} in \autoref{sec:state} is this construction with
$N{=}4$ (base control point $\vR_{k-3}$, three increments).

\subsection{\ac{map} Estimation as Weighted Least Squares}
\label{sec:prelim_map}
The control points and biases above form the finite-dimensional state
$\calX$ (\autoref{sec:problem}).  Each measurement $\vz_i$ relates to
the state through a known model $\mathbf{h}_i$,
$\vz_i = \mathbf{h}_i(\calX) + \mathbf{n}_i$, with the noise
$\mathbf{n}_i$ zero-mean Gaussian of covariance
$\boldsymbol{\Sigma}_i$ and independent across measurements.  The
likelihood is then
$p(\vz\mid\calX)\propto\prod_i
 \exp(-\tfrac12\mathbf{r}_i^{\top}\boldsymbol{\Sigma}_i^{-1}\mathbf{r}_i)$
with residuals $\mathbf{r}_i=\vz_i-\mathbf{h}_i(\calX)$, so its negative
logarithm is a weighted sum of squared residuals: maximum-likelihood
estimation minimizes
$\tfrac12\sum_i\mathbf{r}_i^{\top}\boldsymbol{\Sigma}_i^{-1}\mathbf{r}_i$,
in which each residual is weighted by its inverse covariance
(information) $\boldsymbol{\Sigma}_i^{-1}$.  For a scalar residual this
is a scalar weight $\lambda_i=1/\sigma_i^2$: a more certain sensor is
trusted more.  The variance measurements in \autoref{sec:char}
inform the stream weights $\lambda_a,\lambda_g$ of
\autoref{sec:models} and the per-return radar weights of
\autoref{sec:robust}.

Prior knowledge $p(\calX)$, here trajectory smoothness and, online, the
marginalization prior summarizing past data, turns the estimate into a
\ac{map} one, $\arg\max_{\calX} p(\vz\mid\calX)\,p(\calX)$; its negative
log contributes the extra quadratic terms $E_{\text{reg}}$ and
$E_{\text{prior}}$.  In a sliding window the prior is itself a
posterior: the previous window's posterior over the states both
windows share, with the states that left the window integrated out
(\autoref{sec:sw}).  Radar Doppler residuals have heavy tails
(\autoref{sec:char_noise} measures both the Gaussian core and the
tails); the identifiable outliers are removed before the solve
(\autoref{sec:char_outliers}), so every residual that reaches the
cost enters it quadratically.  Collecting all terms gives the weighted nonlinear
least-squares objective instantiated for our sensors in
\eqref{eq:map}.  We minimize it with Levenberg--Marquardt and
analytic Jacobians.  Rotation updates are applied on the manifold,
$\vR\mapsto\vR\,\Exp(\boldsymbol{\delta})$, so every iterate
remains a valid rotation.  The Jacobians use spline values and
derivatives at each measurement's exact timestamp, provided in
closed form by \autoref{sec:prelim_bspline} and
\autoref{sec:prelim_so3}.
